% W. H. Mills, "A prime-representing function",
% Bull. Amer. Math. Soc. 53 (1947), 604, reset in Mills Modern.
% Variants, chosen by defining a macro before % W. H. Mills, "A prime-representing function",
% Bull. Amer. Math. Soc. 53 (1947), 604, reset in Mills Modern.
% Variants, chosen by defining a macro before % W. H. Mills, "A prime-representing function",
% Bull. Amer. Math. Soc. 53 (1947), 604, reset in Mills Modern.
% Variants, chosen by defining a macro before % W. H. Mills, "A prime-representing function",
% Bull. Amer. Math. Soc. 53 (1947), 604, reset in Mills Modern.
% Variants, chosen by defining a macro before \input{mills}:
%   (none)              Mills Modern (pdflatex)
%   \def\usecm{}        plain Computer Modern (pdflatex)
%   \def\useoldstandard{} Old Standard text + math (lualatex)
%   \def\useeighta{}    Mills 8A, the revival traced from 1947 scans (lualatex)
\ifdefined\useeighta
\documentclass[11pt]{article}
\else
\documentclass[10pt]{article}
\fi
\usepackage[leqno]{amsmath}
\ifdefined\useoldstandard
  \usepackage{unicode-math}
  \setmainfont{OldStandard-Regular.otf}[Scale=1.06,
    ItalicFont=OldStandard-Italic.otf,
    BoldFont=OldStandard-Bold.otf,
    BoldItalicFont=OldStandard-BoldItalic.otf]
  \setmathfont{OldStandard-Math.otf}[Scale=1.06]
  \AtBeginDocument{\thickmuskip=3.5mu plus 2mu}
\else\ifdefined\useeighta
  % 11pt Monotype Modern 8A on 12pt, 317pt measure, as in the 1947 Bulletin
  \usepackage{unicode-math}
  \setmainfont{Mills8A-Regular.otf}[
    ItalicFont=Mills8A-Italic.otf,
    BoldFont=Mills8A-Regular.otf, BoldFeatures={FakeBold=2}]
  \setmathfont{latinmodern-math.otf}
  \setmathfont{Mills8A-Italic.otf}[range={it/{latin,Latin},
    "1D6FC,"1D6FD,"1D701,"1D70B,"1D70E,"1D716,"1D719}]
  \setmathfont{Mills8A-Regular.otf}[range={up/num,
    "2B,"3D,"3C,"3E,"28,"29,"5B,"5D,"2F,"7C,"2212,"2266,"2267,"221E,"2192}]
  \makeatletter
  \renewcommand\normalsize{\@setfontsize\normalsize{11}{12}%
    \abovedisplayskip 6pt plus 2pt minus 2pt
    \belowdisplayskip \abovedisplayskip
    \abovedisplayshortskip 3pt plus 2pt
    \belowdisplayshortskip 3pt plus 2pt}
  \renewcommand\footnotesize{\@setfontsize\footnotesize{9}{10}}
  \makeatother
  \AtBeginDocument{\normalsize\thickmuskip=3.5mu plus 2mu}
\else\ifdefined\usecm\else
  \usepackage{millsmodern}
\fi\fi\fi
\ifdefined\useeighta
\usepackage[paperwidth=31pc,paperheight=50pc,textwidth=317pt,
  top=4pc,bottom=4pc,footskip=2pc]{geometry}
\else
\usepackage[paperwidth=31pc,paperheight=50pc,textwidth=26pc,
  top=4pc,bottom=4pc,footskip=2pc]{geometry}
\fi
\ifdefined\pdfpkresolution\pdfpkresolution=1200 \fi
\setlength\parindent{1.5em}
\setlength\parskip{0pt}
\ifdefined\useeighta\else\linespread{1.04}\fi
\pagestyle{empty}
\ExplSyntaxOn
\NewDocumentCommand\spaced{m}{% letterspaced title lines
  \seq_set_split:Nnn \l_tmpa_seq {~} {#1}
  \seq_map_indexed_inline:Nn \l_tmpa_seq
    { \int_compare:nT {##1>1} {\hspace{.45em}}
      \tl_map_inline:nn {##2} {####1\kern.07em} \unkern }}
\ExplSyntaxOff
\newcommand\thmhead[1]{\par\smallskip\begingroup\textsc{#1.}\ \itshape\ignorespaces}
\newcommand\thmend{\par\endgroup\smallskip}
\newcommand\proofhead{\par\textsc{Proof.}\ \ignorespaces}
\renewcommand\thefootnote{\arabic{footnote}}
\makeatletter
\renewcommand\@makefntext[1]{\parindent1.5em\noindent
  \hbox to 1.5em{\hss\@makefnmark\,}#1}
\makeatother
\begin{document}
\begin{center}
  {\bfseries\spaced{A PRIME-REPRESENTING FUNCTION}}\par
  \medskip
  {\footnotesize\spaced{W. H. MILLS}}
\end{center}
\medskip

A function $f(x)$ is said to be a prime-representing function if $f(x)$
is a prime number for all positive integral values of~$x$. It will be
shown that there exists a real number $A$ such that $[A^{3^x}]$ is a
prime-representing function, where $[R]$ denotes the greatest integer
less than or equal to~$R$.

Let $p_n$ denote the $n$th prime number. A. E. Ingham\footnote{A. E.
Ingham, \emph{On the difference between consecutive primes}, Quart.\ J.
Math.\ Oxford Ser.\ vol.~8 (1937) pp.~255--266.} has shown that
\begin{equation}
  p_{n+1} - p_n < K p_n^{5/8}
\end{equation}
where $K$ is a fixed positive integer.

\thmhead{Lemma}
If $N$ is an integer greater than $K^8$ there exists a prime $p$ such
that $N^3 < p < (N+1)^3 - 1$.
\thmend

\proofhead
Let $p_n$ be the greatest prime less than $N^3$. Then
\begin{equation}
\begin{split}
  N^3 &< p_{n+1} < p_n + K p_n^{5/8} < N^3 + K N^{15/8} < N^3 + N^2 \\
      &< (N+1)^3 - 1.
\end{split}
\end{equation}

Let $P_0$ be a prime greater than $K^8$. Then by the lemma we can
construct an infinite sequence of primes, $P_0, P_1, P_2, \cdots$, such
that $P_n^3 < P_{n+1} < (P_n+1)^3 - 1$. Let
\begin{equation}
  u_n = P_n^{3^{-n}}, \qquad v_n = (P_n + 1)^{3^{-n}}.
\end{equation}
Then
\begin{gather}
  v_n > u_n, \qquad u_{n+1} = P_{n+1}^{3^{-n-1}} > P_n^{3^{-n}} = u_n, \\
  v_{n+1} = (P_{n+1} + 1)^{3^{-n-1}} < (P_n + 1)^{3^{-n}} = v_n.
\end{gather}
It follows at once that the $u_n$ form a bounded monotone increasing
sequence. Let $A = \lim_{n\to\infty} u_n$.

\thmhead{Theorem}
$[A^{3^x}]$ is a prime-representing function.
\thmend

\proofhead
From (4) and (5) it follows that $u_n < A < v_n$, or
$P_n < A^{3^n} < P_n + 1$. Therefore $[A^{3^n}] = P_n$ and $[A^{3^x}]$
is a prime-\hspace{0pt}representing function.
\end{document}
:
%   (none)              Mills Modern (pdflatex)
%   \def\usecm{}        plain Computer Modern (pdflatex)
%   \def\useoldstandard{} Old Standard text + math (lualatex)
%   \def\useeighta{}    Mills 8A, the revival traced from 1947 scans (lualatex)
\ifdefined\useeighta
\documentclass[11pt]{article}
\else
\documentclass[10pt]{article}
\fi
\usepackage[leqno]{amsmath}
\ifdefined\useoldstandard
  \usepackage{unicode-math}
  \setmainfont{OldStandard-Regular.otf}[Scale=1.06,
    ItalicFont=OldStandard-Italic.otf,
    BoldFont=OldStandard-Bold.otf,
    BoldItalicFont=OldStandard-BoldItalic.otf]
  \setmathfont{OldStandard-Math.otf}[Scale=1.06]
  \AtBeginDocument{\thickmuskip=3.5mu plus 2mu}
\else\ifdefined\useeighta
  % 11pt Monotype Modern 8A on 12pt, 317pt measure, as in the 1947 Bulletin
  \usepackage{unicode-math}
  \setmainfont{Mills8A-Regular.otf}[
    ItalicFont=Mills8A-Italic.otf,
    BoldFont=Mills8A-Regular.otf, BoldFeatures={FakeBold=2}]
  \setmathfont{latinmodern-math.otf}
  \setmathfont{Mills8A-Italic.otf}[range={it/{latin,Latin},
    "1D6FC,"1D6FD,"1D701,"1D70B,"1D70E,"1D716,"1D719}]
  \setmathfont{Mills8A-Regular.otf}[range={up/num,
    "2B,"3D,"3C,"3E,"28,"29,"5B,"5D,"2F,"7C,"2212,"2266,"2267,"221E,"2192}]
  \makeatletter
  \renewcommand\normalsize{\@setfontsize\normalsize{11}{12}%
    \abovedisplayskip 6pt plus 2pt minus 2pt
    \belowdisplayskip \abovedisplayskip
    \abovedisplayshortskip 3pt plus 2pt
    \belowdisplayshortskip 3pt plus 2pt}
  \renewcommand\footnotesize{\@setfontsize\footnotesize{9}{10}}
  \makeatother
  \AtBeginDocument{\normalsize\thickmuskip=3.5mu plus 2mu}
\else\ifdefined\usecm\else
  \usepackage{millsmodern}
\fi\fi\fi
\ifdefined\useeighta
\usepackage[paperwidth=31pc,paperheight=50pc,textwidth=317pt,
  top=4pc,bottom=4pc,footskip=2pc]{geometry}
\else
\usepackage[paperwidth=31pc,paperheight=50pc,textwidth=26pc,
  top=4pc,bottom=4pc,footskip=2pc]{geometry}
\fi
\ifdefined\pdfpkresolution\pdfpkresolution=1200 \fi
\setlength\parindent{1.5em}
\setlength\parskip{0pt}
\ifdefined\useeighta\else\linespread{1.04}\fi
\pagestyle{empty}
\ExplSyntaxOn
\NewDocumentCommand\spaced{m}{% letterspaced title lines
  \seq_set_split:Nnn \l_tmpa_seq {~} {#1}
  \seq_map_indexed_inline:Nn \l_tmpa_seq
    { \int_compare:nT {##1>1} {\hspace{.45em}}
      \tl_map_inline:nn {##2} {####1\kern.07em} \unkern }}
\ExplSyntaxOff
\newcommand\thmhead[1]{\par\smallskip\begingroup\textsc{#1.}\ \itshape\ignorespaces}
\newcommand\thmend{\par\endgroup\smallskip}
\newcommand\proofhead{\par\textsc{Proof.}\ \ignorespaces}
\renewcommand\thefootnote{\arabic{footnote}}
\makeatletter
\renewcommand\@makefntext[1]{\parindent1.5em\noindent
  \hbox to 1.5em{\hss\@makefnmark\,}#1}
\makeatother
\begin{document}
\begin{center}
  {\bfseries\spaced{A PRIME-REPRESENTING FUNCTION}}\par
  \medskip
  {\footnotesize\spaced{W. H. MILLS}}
\end{center}
\medskip

A function $f(x)$ is said to be a prime-representing function if $f(x)$
is a prime number for all positive integral values of~$x$. It will be
shown that there exists a real number $A$ such that $[A^{3^x}]$ is a
prime-representing function, where $[R]$ denotes the greatest integer
less than or equal to~$R$.

Let $p_n$ denote the $n$th prime number. A. E. Ingham\footnote{A. E.
Ingham, \emph{On the difference between consecutive primes}, Quart.\ J.
Math.\ Oxford Ser.\ vol.~8 (1937) pp.~255--266.} has shown that
\begin{equation}
  p_{n+1} - p_n < K p_n^{5/8}
\end{equation}
where $K$ is a fixed positive integer.

\thmhead{Lemma}
If $N$ is an integer greater than $K^8$ there exists a prime $p$ such
that $N^3 < p < (N+1)^3 - 1$.
\thmend

\proofhead
Let $p_n$ be the greatest prime less than $N^3$. Then
\begin{equation}
\begin{split}
  N^3 &< p_{n+1} < p_n + K p_n^{5/8} < N^3 + K N^{15/8} < N^3 + N^2 \\
      &< (N+1)^3 - 1.
\end{split}
\end{equation}

Let $P_0$ be a prime greater than $K^8$. Then by the lemma we can
construct an infinite sequence of primes, $P_0, P_1, P_2, \cdots$, such
that $P_n^3 < P_{n+1} < (P_n+1)^3 - 1$. Let
\begin{equation}
  u_n = P_n^{3^{-n}}, \qquad v_n = (P_n + 1)^{3^{-n}}.
\end{equation}
Then
\begin{gather}
  v_n > u_n, \qquad u_{n+1} = P_{n+1}^{3^{-n-1}} > P_n^{3^{-n}} = u_n, \\
  v_{n+1} = (P_{n+1} + 1)^{3^{-n-1}} < (P_n + 1)^{3^{-n}} = v_n.
\end{gather}
It follows at once that the $u_n$ form a bounded monotone increasing
sequence. Let $A = \lim_{n\to\infty} u_n$.

\thmhead{Theorem}
$[A^{3^x}]$ is a prime-representing function.
\thmend

\proofhead
From (4) and (5) it follows that $u_n < A < v_n$, or
$P_n < A^{3^n} < P_n + 1$. Therefore $[A^{3^n}] = P_n$ and $[A^{3^x}]$
is a prime-\hspace{0pt}representing function.
\end{document}
:
%   (none)              Mills Modern (pdflatex)
%   \def\usecm{}        plain Computer Modern (pdflatex)
%   \def\useoldstandard{} Old Standard text + math (lualatex)
%   \def\useeighta{}    Mills 8A, the revival traced from 1947 scans (lualatex)
\ifdefined\useeighta
\documentclass[11pt]{article}
\else
\documentclass[10pt]{article}
\fi
\usepackage[leqno]{amsmath}
\ifdefined\useoldstandard
  \usepackage{unicode-math}
  \setmainfont{OldStandard-Regular.otf}[Scale=1.06,
    ItalicFont=OldStandard-Italic.otf,
    BoldFont=OldStandard-Bold.otf,
    BoldItalicFont=OldStandard-BoldItalic.otf]
  \setmathfont{OldStandard-Math.otf}[Scale=1.06]
  \AtBeginDocument{\thickmuskip=3.5mu plus 2mu}
\else\ifdefined\useeighta
  % 11pt Monotype Modern 8A on 12pt, 317pt measure, as in the 1947 Bulletin
  \usepackage{unicode-math}
  \setmainfont{Mills8A-Regular.otf}[
    ItalicFont=Mills8A-Italic.otf,
    BoldFont=Mills8A-Regular.otf, BoldFeatures={FakeBold=2}]
  \setmathfont{latinmodern-math.otf}
  \setmathfont{Mills8A-Italic.otf}[range={it/{latin,Latin},
    "1D6FC,"1D6FD,"1D701,"1D70B,"1D70E,"1D716,"1D719}]
  \setmathfont{Mills8A-Regular.otf}[range={up/num,
    "2B,"3D,"3C,"3E,"28,"29,"5B,"5D,"2F,"7C,"2212,"2266,"2267,"221E,"2192}]
  \makeatletter
  \renewcommand\normalsize{\@setfontsize\normalsize{11}{12}%
    \abovedisplayskip 6pt plus 2pt minus 2pt
    \belowdisplayskip \abovedisplayskip
    \abovedisplayshortskip 3pt plus 2pt
    \belowdisplayshortskip 3pt plus 2pt}
  \renewcommand\footnotesize{\@setfontsize\footnotesize{9}{10}}
  \makeatother
  \AtBeginDocument{\normalsize\thickmuskip=3.5mu plus 2mu}
\else\ifdefined\usecm\else
  \usepackage{millsmodern}
\fi\fi\fi
\ifdefined\useeighta
\usepackage[paperwidth=31pc,paperheight=50pc,textwidth=317pt,
  top=4pc,bottom=4pc,footskip=2pc]{geometry}
\else
\usepackage[paperwidth=31pc,paperheight=50pc,textwidth=26pc,
  top=4pc,bottom=4pc,footskip=2pc]{geometry}
\fi
\ifdefined\pdfpkresolution\pdfpkresolution=1200 \fi
\setlength\parindent{1.5em}
\setlength\parskip{0pt}
\ifdefined\useeighta\else\linespread{1.04}\fi
\pagestyle{empty}
\ExplSyntaxOn
\NewDocumentCommand\spaced{m}{% letterspaced title lines
  \seq_set_split:Nnn \l_tmpa_seq {~} {#1}
  \seq_map_indexed_inline:Nn \l_tmpa_seq
    { \int_compare:nT {##1>1} {\hspace{.45em}}
      \tl_map_inline:nn {##2} {####1\kern.07em} \unkern }}
\ExplSyntaxOff
\newcommand\thmhead[1]{\par\smallskip\begingroup\textsc{#1.}\ \itshape\ignorespaces}
\newcommand\thmend{\par\endgroup\smallskip}
\newcommand\proofhead{\par\textsc{Proof.}\ \ignorespaces}
\renewcommand\thefootnote{\arabic{footnote}}
\makeatletter
\renewcommand\@makefntext[1]{\parindent1.5em\noindent
  \hbox to 1.5em{\hss\@makefnmark\,}#1}
\makeatother
\begin{document}
\begin{center}
  {\bfseries\spaced{A PRIME-REPRESENTING FUNCTION}}\par
  \medskip
  {\footnotesize\spaced{W. H. MILLS}}
\end{center}
\medskip

A function $f(x)$ is said to be a prime-representing function if $f(x)$
is a prime number for all positive integral values of~$x$. It will be
shown that there exists a real number $A$ such that $[A^{3^x}]$ is a
prime-representing function, where $[R]$ denotes the greatest integer
less than or equal to~$R$.

Let $p_n$ denote the $n$th prime number. A. E. Ingham\footnote{A. E.
Ingham, \emph{On the difference between consecutive primes}, Quart.\ J.
Math.\ Oxford Ser.\ vol.~8 (1937) pp.~255--266.} has shown that
\begin{equation}
  p_{n+1} - p_n < K p_n^{5/8}
\end{equation}
where $K$ is a fixed positive integer.

\thmhead{Lemma}
If $N$ is an integer greater than $K^8$ there exists a prime $p$ such
that $N^3 < p < (N+1)^3 - 1$.
\thmend

\proofhead
Let $p_n$ be the greatest prime less than $N^3$. Then
\begin{equation}
\begin{split}
  N^3 &< p_{n+1} < p_n + K p_n^{5/8} < N^3 + K N^{15/8} < N^3 + N^2 \\
      &< (N+1)^3 - 1.
\end{split}
\end{equation}

Let $P_0$ be a prime greater than $K^8$. Then by the lemma we can
construct an infinite sequence of primes, $P_0, P_1, P_2, \cdots$, such
that $P_n^3 < P_{n+1} < (P_n+1)^3 - 1$. Let
\begin{equation}
  u_n = P_n^{3^{-n}}, \qquad v_n = (P_n + 1)^{3^{-n}}.
\end{equation}
Then
\begin{gather}
  v_n > u_n, \qquad u_{n+1} = P_{n+1}^{3^{-n-1}} > P_n^{3^{-n}} = u_n, \\
  v_{n+1} = (P_{n+1} + 1)^{3^{-n-1}} < (P_n + 1)^{3^{-n}} = v_n.
\end{gather}
It follows at once that the $u_n$ form a bounded monotone increasing
sequence. Let $A = \lim_{n\to\infty} u_n$.

\thmhead{Theorem}
$[A^{3^x}]$ is a prime-representing function.
\thmend

\proofhead
From (4) and (5) it follows that $u_n < A < v_n$, or
$P_n < A^{3^n} < P_n + 1$. Therefore $[A^{3^n}] = P_n$ and $[A^{3^x}]$
is a prime-\hspace{0pt}representing function.
\end{document}
:
%   (none)              Mills Modern (pdflatex)
%   \def\usecm{}        plain Computer Modern (pdflatex)
%   \def\useoldstandard{} Old Standard text + math (lualatex)
%   \def\useeighta{}    Mills 8A, the revival traced from 1947 scans (lualatex)
\ifdefined\useeighta
\documentclass[11pt]{article}
\else
\documentclass[10pt]{article}
\fi
\usepackage[leqno]{amsmath}
\ifdefined\useoldstandard
  \usepackage{unicode-math}
  \setmainfont{OldStandard-Regular.otf}[Scale=1.06,
    ItalicFont=OldStandard-Italic.otf,
    BoldFont=OldStandard-Bold.otf,
    BoldItalicFont=OldStandard-BoldItalic.otf]
  \setmathfont{OldStandard-Math.otf}[Scale=1.06]
  \AtBeginDocument{\thickmuskip=3.5mu plus 2mu}
\else\ifdefined\useeighta
  % 11pt Monotype Modern 8A on 12pt, 317pt measure, as in the 1947 Bulletin
  \usepackage{unicode-math}
  \setmainfont{Mills8A-Regular.otf}[
    ItalicFont=Mills8A-Italic.otf,
    BoldFont=Mills8A-Regular.otf, BoldFeatures={FakeBold=2}]
  \setmathfont{latinmodern-math.otf}
  \setmathfont{Mills8A-Italic.otf}[range={it/{latin,Latin},
    "1D6FC,"1D6FD,"1D701,"1D70B,"1D70E,"1D716,"1D719}]
  \setmathfont{Mills8A-Regular.otf}[range={up/num,
    "2B,"3D,"3C,"3E,"28,"29,"5B,"5D,"2F,"7C,"2212,"2266,"2267,"221E,"2192}]
  \makeatletter
  \renewcommand\normalsize{\@setfontsize\normalsize{11}{12}%
    \abovedisplayskip 6pt plus 2pt minus 2pt
    \belowdisplayskip \abovedisplayskip
    \abovedisplayshortskip 3pt plus 2pt
    \belowdisplayshortskip 3pt plus 2pt}
  \renewcommand\footnotesize{\@setfontsize\footnotesize{9}{10}}
  \makeatother
  \AtBeginDocument{\normalsize\thickmuskip=3.5mu plus 2mu}
\else\ifdefined\usecm\else
  \usepackage{millsmodern}
\fi\fi\fi
\ifdefined\useeighta
\usepackage[paperwidth=31pc,paperheight=50pc,textwidth=317pt,
  top=4pc,bottom=4pc,footskip=2pc]{geometry}
\else
\usepackage[paperwidth=31pc,paperheight=50pc,textwidth=26pc,
  top=4pc,bottom=4pc,footskip=2pc]{geometry}
\fi
\ifdefined\pdfpkresolution\pdfpkresolution=1200 \fi
\setlength\parindent{1.5em}
\setlength\parskip{0pt}
\ifdefined\useeighta\else\linespread{1.04}\fi
\pagestyle{empty}
\ExplSyntaxOn
\NewDocumentCommand\spaced{m}{% letterspaced title lines
  \seq_set_split:Nnn \l_tmpa_seq {~} {#1}
  \seq_map_indexed_inline:Nn \l_tmpa_seq
    { \int_compare:nT {##1>1} {\hspace{.45em}}
      \tl_map_inline:nn {##2} {####1\kern.07em} \unkern }}
\ExplSyntaxOff
\newcommand\thmhead[1]{\par\smallskip\begingroup\textsc{#1.}\ \itshape\ignorespaces}
\newcommand\thmend{\par\endgroup\smallskip}
\newcommand\proofhead{\par\textsc{Proof.}\ \ignorespaces}
\renewcommand\thefootnote{\arabic{footnote}}
\makeatletter
\renewcommand\@makefntext[1]{\parindent1.5em\noindent
  \hbox to 1.5em{\hss\@makefnmark\,}#1}
\makeatother
\begin{document}
\begin{center}
  {\bfseries\spaced{A PRIME-REPRESENTING FUNCTION}}\par
  \medskip
  {\footnotesize\spaced{W. H. MILLS}}
\end{center}
\medskip

A function $f(x)$ is said to be a prime-representing function if $f(x)$
is a prime number for all positive integral values of~$x$. It will be
shown that there exists a real number $A$ such that $[A^{3^x}]$ is a
prime-representing function, where $[R]$ denotes the greatest integer
less than or equal to~$R$.

Let $p_n$ denote the $n$th prime number. A. E. Ingham\footnote{A. E.
Ingham, \emph{On the difference between consecutive primes}, Quart.\ J.
Math.\ Oxford Ser.\ vol.~8 (1937) pp.~255--266.} has shown that
\begin{equation}
  p_{n+1} - p_n < K p_n^{5/8}
\end{equation}
where $K$ is a fixed positive integer.

\thmhead{Lemma}
If $N$ is an integer greater than $K^8$ there exists a prime $p$ such
that $N^3 < p < (N+1)^3 - 1$.
\thmend

\proofhead
Let $p_n$ be the greatest prime less than $N^3$. Then
\begin{equation}
\begin{split}
  N^3 &< p_{n+1} < p_n + K p_n^{5/8} < N^3 + K N^{15/8} < N^3 + N^2 \\
      &< (N+1)^3 - 1.
\end{split}
\end{equation}

Let $P_0$ be a prime greater than $K^8$. Then by the lemma we can
construct an infinite sequence of primes, $P_0, P_1, P_2, \cdots$, such
that $P_n^3 < P_{n+1} < (P_n+1)^3 - 1$. Let
\begin{equation}
  u_n = P_n^{3^{-n}}, \qquad v_n = (P_n + 1)^{3^{-n}}.
\end{equation}
Then
\begin{gather}
  v_n > u_n, \qquad u_{n+1} = P_{n+1}^{3^{-n-1}} > P_n^{3^{-n}} = u_n, \\
  v_{n+1} = (P_{n+1} + 1)^{3^{-n-1}} < (P_n + 1)^{3^{-n}} = v_n.
\end{gather}
It follows at once that the $u_n$ form a bounded monotone increasing
sequence. Let $A = \lim_{n\to\infty} u_n$.

\thmhead{Theorem}
$[A^{3^x}]$ is a prime-representing function.
\thmend

\proofhead
From (4) and (5) it follows that $u_n < A < v_n$, or
$P_n < A^{3^n} < P_n + 1$. Therefore $[A^{3^n}] = P_n$ and $[A^{3^x}]$
is a prime-\hspace{0pt}representing function.
\end{document}
